\section{Prólogo}
É evidente que a ciência da robótica depende de conceitos plenamente compreendidos 
mediante investigação matemática e lógica - pois toda máquina opera segundo quantidade, 
relação, movimento e ordem. Com efeito, aquilo que custa a ser provado através deste rigor é incapaz de sustentar-se
de modo cientificamente satisfatório. 
\par
§1 --- De acordo com o \textsc{Matemático}\footnote{Euclides de Alexandria}, é possível 
representar o mundo visível a partir de axiomas lógicos e geométricos. É evidente que o 
quadrado da hipotenusa de um triângulo é a soma dos quadrados dos catetos pois o 
procedimento é demonstrável e observável.
\par
§2 --- Nas máquinas, a observação ocorre de modo distinto. A matemática precisa ser 
utilizada de forma contínua, visto que durante a execução ocorre movimento e cadência 
em tempo real. 
\par
§3 --- Por definição, robôs não são meramente mecânismos físicos. Suas operações são 
estabelecidas através de uma ordem matemática e lógica que deve corresponder às 
alterações de seu estado físico.
\par
Portanto, o estudo da robótica torna-se dependente de um ordenamento matemático para 
suas operações, isto é, correspondência entre a definição e o comportamento físico.
Contudo, este livro destina-se essencialmente a ser uma breve introdução, mas não tende 
a ser uma 'summa scientiae' no que diz a respeito do tópico, apesar de apresentar os 
fundamentos para isto. 